\title{xLSTM: Extended Long Short-Term Memory}

\begin{abstract}
The introduction of the constant error carousel and gating mechanisms during the 1990s laid the groundwork for the Long Short-Term Memory (LSTM) model. Over the years, LSTMs have consistently demonstrated their value, playing pivotal roles in several breakthroughs in deep learning, notably serving as the foundation for the earliest Large Language Models (LLMs). Nevertheless, with the emergence of Transformer architectures—characterized by scalable, parallelizable self-attention—a new standard was set, with LSTMs being surpassed when it came to scaling up. This prompts a fundamental investigation: what level of performance can be attained in language modeling by scaling LSTMs up to the billion-parameter range, applying the latest advances from modern LLMs, and overcoming established weaknesses of conventional LSTMs? To this end, we introduce exponential gating equipped with advanced normalization and stabilization strategies. We also reimagine the memory design in LSTMs, resulting in: (i) sLSTM, which features scalar memory cells, scalar updates, and an innovative mixing approach to memory, and (ii) mLSTM, constructed with fully parallelizable matrix memory and an update mechanism based on covariance. These proposed LSTM modifications, when incorporated into the framework of residual block structures, constitute what we term as xLSTM blocks. Stacking these blocks using residual connections produces comprehensive xLSTM architectures. The synergetic effects of exponential gating and updated memory frameworks significantly enhance the capabilities of xLSTM, allowing these models to compete impressively with cutting-edge Transformers and State Space Models, regarding both efficacy and scalability.\\
Code available at: \url{https://github.com/NX-AI/xlstm}
\end{abstract}

\section{Introduction}
\label{sec:intro}

The Long Short-Term Memory (LSTM) ideas~\citep{Hochreiter:91,Hochreiter:97nips,Hochreiter:97}, i.e., the constant error carousel and gating, were introduced to overcome the vanishing gradient problem of recurrent neural networks~\citep{Hochreiter:91,Hochreiter:00book}:

\begin{align}
\label{eq:lstm_idea}
  c_t \ = \ f_t \ c_{t-1} \ + \ 
          i_t \  z_t \ , \quad  
          h_t \ = \ o_t \ \psi({c_t}) \ . 
\end{align}

The constant error carousel refers to the cell state's $c_{t-1}$ (illustrated in green) being incrementally updated through additive contributions from cell inputs $z_t$, regulated by sigmoid gates (depicted in blue). The input gate $i_t$ alongside the forget gate $f_t$ manage how these updates occur, whereas the output gate $o_t$ determines the release of information from the memory cell, yielding the hidden state $h_t$. Before being gated for output, the cell state undergoes a normalization or squashing transformation via $\psi$, after which the output gate produces the hidden state.

LSTMs have achieved notable success across a wide range of application areas \citep{Hochreiter:01,Hochreiter:07,Schmidhuber:15}, maintaining dominance in text generation tasks prior to the introduction of Transformers in 2017~\citep{Vaswani:17}. Their capability has been validated in an array of sequence-based problems, such as text generation~\citep{Graves:13,Karpathy:15blog}, handwriting synthesis~\citep{Graves:13}, sequence-to-sequence machine translation~\citep{Sutskever:14nips}, program evaluation~\citep{Zaremba:14arxiv}, automatic image captioning~\citep{Karpathy:15,Hossain:19}, code generation~\citep{Karpathy:15blog}, as well as hydrological applications like rainfall-runoff prediction~\citep{Kratzert:18,Kratzert:19} and flood early-warning systems~\citep{Nearing:24}. Within reinforcement learning, LSTM-based sequence models have yielded top performance; prominent examples include the AlphaStar system for StarCraft II~\citep{Vinyals:17short}, OpenAI Five for Dota 2~\citep{OpenAI:19}, and the magnetic control frameworks employed in nuclear fusion reactors~\citep{Degrave:22short}. LSTMs exhibit a marked proficiency for representation learning by effectively extracting and retaining semantic information through their memory states~\citep{Karpathy:15blog}, as demonstrated by the emergence of neurons capturing numerical and syntactic structure~\citep{Lakretz:19}, linguistic properties~\citep{Bau:19}, and sentiment~\citep{Radford:17arxiv}. Despite recent advancements in sequence modeling, LSTMs continue to play a pivotal role in key practical applications~\citep{Degrave:22short,Nearing:24}, underscoring their enduring impact.

Although LSTMs have achieved significant success, they suffer from three principal drawbacks: (i) Inability to modify storage decisions retrospectively. To illustrate this weakness, we consider the {\it Nearest Neighbor Search} problem (refer to Appendix~\ref{sec:appNearestNeighborSearch} for more detail): Given a reference vector, a sequence must be traversed sequentially to locate the most similar vector and retrieve its corresponding value at the sequence’s conclusion. The left panel of Figure~\ref{fig:lstmProblems} presents the mean squared error for this task. LSTMs face difficulty updating a previously stored value when a more similar match arises later in the sequence. Our proposed xLSTM overcomes this challenge by employing exponential gating. (ii) Constrained memory capacity, in that all information is compressed into scalar cell states. This shortcoming is demonstrated with the {\it Rare Token Prediction} task. In the right panel of Figure~\ref{fig:lstmProblems}, the perplexity on Wikitext-103~\citep{Merity:17} is shown for token subsets divided by frequency. LSTMs underperform on rare tokens due to their restricted memory, whereas our new xLSTM leverages a matrix memory to address this. (iii) Inherent lack of parallelism as a result of memory mixing—that is, the recurrent dependencies among hidden states from one timestep to the next require strictly sequential computation. 

These shortcomings of LSTM architectures have driven the adoption of Transformers~\citep{Vaswani:17} in the field of language modeling. This raises an intriguing question: to what extent can language modeling performance be improved by resolving these limitations and scaling LSTM variants to sizes comparable to contemporary Large Language Models?

\section{Extended Long Short-Term Memory}
\label{sec:xLSTM}

In response to the limitations found in LSTM, Extended Long Short-Term Memory (xLSTM) presents two significant advances over the conventional LSTM formulation outlined in Equation~\eqref{eq:lstm_idea}. The first advance involves the introduction of exponential gating, while the second concerns the adoption of novel memory structures. These enhancements result in two distinctive extensions of the LSTM framework: (i) the sLSTM (see Section~\ref{sec:sLSTM}), characterized by a scalar memory and update mechanism with memory mixing capabilities, and (ii) the mLSTM (see Section~\ref{sec:mLSTM}), which utilizes a matrix-based memory and a covariance-based update scheme (using outer products) that fully supports parallel computation. Exponential gating underpins improvements in both the sLSTM and mLSTM models. In particular, the mLSTM achieves parallelism by removing memory mixing and, consequently, hidden-to-hidden recurrence. Both variants, sLSTM and mLSTM, are adaptable to configurations with several memory cells; in these settings, sLSTM continues to mix information across cells. Additionally, sLSTM can feature multiple independent heads—each head conducting memory mixing among its own internal cells, but not across other heads—thus introducing a new strategy for information fusion alongside exponential gating. In the case of mLSTM, the concepts of multiple heads and multiple cells coincide.

When these novel LSTM variants are incorporated into residual block architectures, they form xLSTM blocks (see Section~\ref{sec:xLSTMblock}). Constructing neural networks by residually stacking xLSTM blocks yields xLSTM architectures, elaborated in Section~\ref{sec:xLSTMarchitecure}. A schematic of this architecture and its structural elements appears in Figure~\ref{fig:xlstm_sketch}.

\subsection{Review of the Long Short-Term Memory}
\label{sec:orgLSTM}

The foundational LSTM framework~\citep{Hochreiter:91,Hochreiter:97nips,Hochreiter:97} proposed the scalar memory cell as its core mechanism for computation and information retention. This design addresses the vanishing gradient problem~\citep{Hochreiter:91,Hochreiter:00book} via the so-called constant error carousel, which ensures the persistence of gradient flow during cell state updates. Three types of gates regulate the memory cell: the input gate, output gate, and, subsequently introduced by \citet{Gers:00}, the forget gate. The LSTM cell update equations for a given time step $t$ are as follows:

\begin{align}
c_t \ &= \  f_t \ c_{t-1} \ + \ i_t \ z_t &  & &\text{cell state} \\
h_t  \ &= \ o_t \ \tilde{h}_t \ , & \tilde{h}_t \ &= \ \psi \left( c_t\right)
  &\text{hidden state} \\
z_t \ &= \ \varphi \left( \tilde{z}_t \right) \ , 
  &\tilde{z}_t \ &=  \ \mathbf{w}^\top_{z} \ \mathbf{x}_t \ + \
  r_{z}  h_{t-1} \ + \  b_{z} \ \
  &\text{cell input} \\
i_t \ &= \ \sigma \left( \tilde{i}_t  \right) \ , 
  &\tilde{i}_t \ &= \ \mathbf{w}^\top_{i} \ \mathbf{x}_t \ + \
  r_{i}  \ h_{t-1} \ + \  b_{i} \ \
  &\text{input gate} \\
f_t \ &= \ \sigma \left(  \tilde{f}_t \right) \ , 
  &\tilde{f}_t \ &= \ \mathbf{w}^\top_{f} \ \mathbf{x}_t  \ + \
  r_{f}  \ h_{t-1} \ + \  b_{f} \ \
  &\text{forget gate} \\
o_t \ &= \ \sigma \left( \tilde{o}_t \right) \ , 
  &\tilde{o}_t  \ &= \ \mathbf{w}^\top_{o} \ \mathbf{x}_t \ + \
  r_{o}  \ h_{t-1} \ + \  b_{o} \ \
  &\text{output gate} 
\end{align}

The vectors $\mathbf{w}_{z}$, $\mathbf{w}_{i}$, $\mathbf{w}_{f}$, and $\mathbf{w}_{o}$ represent the weights linking the input $\mathbf{x}_t$ to the cell input, input gate, forget gate, and output gate, respectively. Similarly, $r_{z}$, $r_{i}$, $r_{f}$, and $r_{o}$ denote the recurrent weights from the previous hidden state $h_{t-1}$ to the cell input, input gate, forget gate, and output gate. Biases associated with each component are given by $b_{z}$, $b_{i}$, $b_{f}$, and $b_{o}$. The activation functions $\varphi$ and $\psi$ pertain to the cell input and hidden state, respectively, and are most often chosen as $\tanh$. The function $\psi$ ensures the cell state remains bounded by normalizing or compressing it. The activation function for the gates employs the sigmoid function, $\sigma \left( x \right)= 1/(1+\exp(-x))$. Later models incorporated multiple scalar memory cells, $\mathbf{c}_t \in \mathbb{R}^{d}$, within a vector form, facilitating the use of recurrent matrices $\mathbf{R} \in \mathbb{R}^{d \times d}$ to enable the interaction among outputs from different memory cells~\citep{Greff:15}; further elaboration is available in Appendix~\ref{sec:apporgLSTM}. Investigations through ablation confirmed the necessity of each memory cell component~\citep{Greff:15}.

\subsection{sLSTM}
\label{sec:sLSTM}

In order to enable LSTMs to modify their memory choices, we propose the integration of exponential gates (highlighted in red) along with mechanisms for normalization and stabilization. Specifically, we allow the input and forget gates to utilize exponential activation functions. For the purpose of normalization, we add a dedicated normalizer state designed to accumulate the sum of products formed by the input gate and the corresponding future forget gates.

The scalar sLSTM forward pass is:

\begin{align}
\label{eq:slstmforward}
c_t \ &= \  f_t \ c_{t-1} \ + \ i_t \ z_t & & 
 &\text{cell state} \\
n_t \ &= \  f_t \ n_{t-1} \ + \ i_t  & & 
 &\text{normalizer state} \\
h_t \ &= \ o_t \ \tilde{h}_t \ ,
  &\tilde{h}_t \ &= \ c_t / n_t
&\text{hidden state} \\
z_t \ &= \ \varphi \left( \tilde{z}_t \right) \ , 
  &\tilde{z}_t \ &=  \ \mathbf{w}^\top_{z} \ \mathbf{x}_t \ + \
  r_{z}  h_{t-1} \ + \  b_{z}
  &\text{cell input} \\
i_t \ &= \ \!\exp \left( \tilde{i}_t  \right) \ , 
  &\tilde{i}_t \ &= \ \mathbf{w}^\top_{i} \ \mathbf{x}_t \ + \
  r_{i}  \ h_{t-1} \ + \  b_{i}
  &\text{input gate} \\
f_t \ &= \ \sigma \left(  \tilde{f}_t \right) \ \text{OR} \
\!\exp \left(  \tilde{f}_t \right) \ ,
  &\tilde{f}_t \ &= \ \mathbf{w}^\top_{f} \ \mathbf{x}_t  \ + \
  r_{f}  \ h_{t-1} \ + \  b_{f}
  &\text{forget gate} \\
o_t \ &= \ \sigma \left( \tilde{o}_t \right) \ , 
  &\tilde{o}_t  \ &= \ \mathbf{w}^\top_{o} \ \mathbf{x}_t \ + \
  r_{o}  \ h_{t-1} \ + \  b_{o}
  &\text{output gate} 
\end{align}

We transfer the original LSTM gating techniques, i.e., input- and/or hidden-dependent gating plus bias term, to the new architectures. Exponential activation functions can lead to large values that cause overflows. Therefore, we stabilize gates with an additional state \mbox{$m_t$~\citep{Milakov:18arxiv}}:

\begin{align}
\label{eq:slstmstabil}
m_t \ &= \ \max \left( \log ( f_t ) + m_{t-1} , \log ( i_t ) \right) &\text{stabilizer state} \\
i'_t \ &= \ \!\exp \left( \log \left ( i_t \right) - m_t \right) \ = \ \!\exp \left( \tilde{i}_t - m_t \right) \ \, 
  &\text{stabil. input gate} \\
f'_t \ &= \ \!\exp \left( \log \left( f_t \right) + m_{t-1} - m_t \right) \ \, 
  &\text{stabil. forget gate}
\end{align}

As demonstrated in Appendix~\ref{sec:appsLSTM}, substituting $f_t$ with $f'_t$ and $i_t$ with $i'_t$ during the forward computation leaves both the overall network output and the loss gradients with respect to the parameters unaffected.

\paragraph{New Memory Mixing.} Similar to the standard LSTM, sLSTM is capable of utilizing multiple memory cells (refer to Appendix~\ref{sec:appsLSTM}). Having several memory cells allows for memory mixing, facilitated by recurrent connections $\mathbf{R}_{\mathbf{z}}$, $\mathbf{R}_{\mathbf{i}}$, $\mathbf{R}_{\mathbf{f}}$, $\mathbf{R}_{\mathbf{o}}$ that link the hidden state vector $\mathbf{h}$ to the memory cell input $\mathbf{z}$ as well as to the gates $\mathbf{i}$, $\mathbf{f}$, and $\mathbf{o}$, respectively. Exponential gating introduces a novel dynamic to this memory mixing process. In the updated sLSTM architecture, each head may incorporate its own set of memory mixing mechanisms internally, yet mixing does not occur between different heads. The combination of head-based architecture in sLSTM and exponential gating forms a distinctive strategy for memory mixing.

\subsection{mLSTM}
\label{sec:mLSTM}

To increase the storage capabilities of LSTMs, we replace the traditional scalar memory cell $c \in \mathbb{R}$ with a matrix memory cell $\mathbf{C} \in \mathbb{R}^{d \times d}$. This modification enables retrieval to occur through matrix multiplication operations. At each timestep $t$, our goal is to store a key-value pair comprised of vectors $\mathbf{k}_t \in \mathbb{R}^d$ (key) and $\mathbf{v}_t \in \mathbb{R}^d$ (value), adopting nomenclature from the Transformer architecture. Subsequently, at a later timestep $t + \tau$, we intend for the stored value $\mathbf{v}_t$ to be recoverable via a query vector $\mathbf{q}_{t+\tau} \in \mathbb{R}^d$. This paradigm reflects the framework established in Bidirectional Associative Memories (BAMs) \citep{Kohonen:72,Anderson:72,Nakano:72,Anderson:77}.

The covariance update rule~\citep{Sejnowski:77,Dayan:91} for storing a key-value pair is

\begin{align}
\mathbf{C}_t \ &= \ \mathbf{C}_{t-1} \ + \ \mathbf{v}_{t} \ \mathbf{k}_t^\top \ .
\end{align}

We posit the use of layer normalization prior to mapping inputs onto key and value representations, ensuring these vectors have zero mean. According to~\citet{Dayan:91}, the covariance update method achieves optimality for maximizing the separability among recovered binary vectors, which corresponds to maximizing the signal-to-noise ratio. Enhanced separability can be realized when retrieval is restricted to pairwise interactions, accepting a quadratic computational cost similar to that of attention mechanisms~\citep{Krotov:16,Krotov:17,Ramsauer:21}. The covariance update approach also coincides with the principle underlying Fast Weight Programmers~\citep{Schmidhuber:92ncfastweights,Schlag:21}, where a fixed decay rate is applied to $\mathbf{C}_{t-1}$ and a constant learning rate is applied to $\mathbf{v}_{t} \mathbf{k}_t ^\top$~\citep{Ba:16}. In alignment with this idea, we embed the covariance update into the LSTM architecture, identifying the forget gate with the decay rate, the input gate with the learning rate, and the output gate with the scaling of the extracted vector.

For this matrix memory, the normalizer state is the weighted sum of key vectors, where each key vector is weighted by the input gate and all future forget gates. Again, the normalizer state keeps record of the strength of the gates. Since the dot product between query and normalizer state can be close to zero, we use the absolute value of this dot product and lower bound it by a threshold (typically 1.0) as done previously~\citep{Sun:23arxiv}. The mLSTM forward pass is:

\begin{align}
\label{eq:mlstm_recurrent_begin}
\mathbf{C}_t \ &= \  f_t \ \mathbf{C}_{t-1} \ + \ 
  i_t \ \mathbf{v}_t \ \mathbf{k}_t^\top & &  &\text{cell state} \\
\mathbf{n}_t \ &= \  f_t \ \mathbf{n}_{t-1} \ + \ 
  i_t \ \mathbf{k}_t & &  &\text{normalizer state} \\
\mathbf{h}_t  \ &= \ \mathbf{o}_t \ \odot \ \tilde{\mathbf{h}}_t
\ , \qquad \qquad 
\tilde{\mathbf{h}}_t \ = \   \mathbf{C}_t \mathbf{q}_t \ / \ 
  \max \left\{ |\mathbf{n}_t^\top \mathbf{q}_t|, 1 \right\} 
   &&&\text{hidden state} \\
\mathbf{q}_t \ &= \ \mathbf{W}_q \ \mathbf{x}_t \ + \ \mathbf{b}_q  & &  &\text{query input} \\
\mathbf{k}_t \ &= \ \frac{1}{\sqrt{d}} \mathbf{W}_k \ \mathbf{x}_t \ + \ \mathbf{b}_k  & &  &\text{key input} \\
\mathbf{v}_t \ &= \ \mathbf{W}_v \ \mathbf{x}_t \ + \ \mathbf{b}_v  & &  &\text{value input} \\
i_t \ &= \ \!\exp \!  \! \left( \tilde{i}_t  \right) \ , 
 \qquad \qquad \ \ \ \ \,
  \tilde{i}_t \ = \ \mathbf{w}^\top_{i} \ \mathbf{x}_t \ + \  b_{i} \ \ \ 
  &&&\text{input gate} \\
f_t \ &= \ \sigma \!  \left(  \tilde{f}_t \right) \ \text{OR} \ \!\exp \!  \! \left(  \tilde{f}_t \right) , \ \ \: \!
\tilde{f}_t \ = \ \mathbf{w}^\top_{f} \ \mathbf{x}_t  \ + \
  b_{f}
  &&& \text{forget gate} \\
\label{eq:mlstm_recurrent_end}
\mathbf{o}_t \ &= \ \sigma \left( \tilde{\mathbf{o}}_t \right) \ , \qquad \qquad \quad \ \ \ \,
  \tilde{\mathbf{o}}_t  \ = \ \mathbf{W}_{\mathbf{o}} \ \mathbf{x}_t \ + \
  \mathbf{b}_{\mathbf{o}}
  &&&\text{output gate} 
\end{align}

The architecture of mLSTM permits the use of several memory cells, analogous to the conventional LSTM. In the case of mLSTM, employing multiple heads or multiple cells yields an equivalent outcome due to the absence of memory interaction between the cells. To ensure numerical stability of the exponential gating mechanisms in mLSTM, we apply the same stabilization procedures utilized in sLSTM, as outlined in Equation~\ref{eq:slstmstabil}. Owing to the lack of memory interaction within mLSTM, the recurrence relation admits a formulation that can be parallelized. Additional elaboration is provided in Appendix~\ref{sec:appmLSTM}.

\subsection{xLSTM Architecture}
\label{sec:xLSTMarch}

\paragraph{xLSTM Blocks.}
\label{sec:xLSTMblock}
An xLSTM block is designed to produce a nonlinear abstraction of preceding inputs within a high-dimensional feature space, enabling it to more effectively distinguish among different historical sequences or contextual backgrounds. Discriminating between diverse pasts is fundamental for generating accurate predictions of subsequent sequence components, such as forthcoming tokens. Our approach draws on Cover’s Theorem~\citep{Cover:65}, which posits that nonlinearly embedded data in higher-dimensional spaces are more likely to be linearly separable compared to embeddings within the original, lower-dimensional space. We examine two residual block designs: (i) The first is a residual block featuring a post up-projection step, akin to the mechanism in Transformers. This variant produces a nonlinear summary of previous states in the initial feature space, performs a linear up-projection to a higher-dimensional space, applies a nonlinear activation, and subsequently maps the result linearly back to the original dimension; this can be found in the left panel of Figure~\ref{fig:Backbones} and the third column of Figure~\ref{fig:xlstm_sketch}. A comprehensive schematic is available in Figure~\ref{fig:appsLSTM_detailed} in the appendix. (ii) The second architecture involves a residual block that applies the linear up-projection step prior to nonlinear transformation, as observed in State Space Models.

encodes information from previous steps through a nonlinear process in the high-dimensional latent space, followed by a linear mapping that returns this representation to the input space. When an xLSTM block contains an sLSTM component, the post up-projection block is employed. Alternatively, if the xLSTM includes an mLSTM, the pre up-projection block is selected, reflecting the expanded memory in the high-dimensional space. For further clarification, the reader may consult the leftmost panel of Figure~\ref{fig:Backbones}, the third column in Figure~\ref{fig:xlstm_sketch}, or Figure~\ref{fig:appsLSTM_detailed} in the appendix.

\paragraph{xLSTM Architecture.}
\label{sec:xLSTMarchitecure}
The xLSTM architecture is formed by stacking building blocks in a residual configuration~\citep{Srivastava:15,He:16}. We adopt pre-LayerNorm~\citep{Ba:16b} residual architectures, a practice prevalent in current Large Language Models. This configuration is illustrated in the final column of Figure~\ref{fig:xlstm_sketch}.

\subsection{Memory and Speed Considerations}

In contrast to Transformers, xLSTM networks offer linear computational complexity and fixed memory requirements regardless of sequence length. Owing to the compressive nature of xLSTM memory, these models are particularly advantageous for deployment in industrial environments and edge devices.

The mLSTM model leverages a non-parameterized memory, but this comes at the cost of considerable computational demand, necessitating both $d \times d$ memory and $d \times d$ update operations. Here, we balance increased memory capability with higher computational overhead. However, leveraging GPU parallelization means that these computations marginally impact total processing time.

Although mLSTM supports parallel execution akin to FlashAttention~\citep{Dao:22,Dao:23} and GLA~\citep{Yang:23arxiv}, sLSTM’s memory interleaving (i.e., hidden-hidden interactions) precludes such parallelization. Nevertheless, we have engineered an efficient CUDA-based implementation that applies memory optimizations down to the register level, resulting in a runtime generally less than twice as long as that of mLSTM.

\section{Related Work}
\label{sec:related}

\paragraph{Linear Attention.} Numerous approaches have been proposed to address the Transformer’s quadratic computational cost with respect to context length, seeking to render attention mechanisms linear. The Synthesizer replaces traditional token-to-token attention with synthetic attention weights learned without explicit token interactions~\citep{Tay:20arxiv}. Linformer achieves a linear self-attention formulation via projecting onto a low-rank matrix, effectively providing a linear approximation~\citep{Wang:20arxiv}. The Linear Transformer reformulates the attention operation to operate in linear time \citep{Katharopoulos:20}. Performer approximates the softmax attention mechanism through a positive orthogonal random features technique~\citep{Choromanski:21}. Additionally, fast long-range convolutions have been adopted as an alternative to attention in models such as Structured Global Convolution (SGConv)~\citep{Li:22} and Hyena Hierarchy~\citep{Poli:23}.

\paragraph{State Space Models.} In recent years, State Space Models (SSMs) have gained considerable traction due to their linear complexity with respect to sequence length and their strong performance as alternatives to Transformers. Early developments in this area include the introduction of the Structured State Space sequence model (S4)~\citep{Gu:21}. This was soon followed by models such as the Diagonal State Space (DSS) model~\citep{Gupta:22}, Gated State Space (GSS) models~\citep{Mehta:22}, the S5 model~\citep{Smith:22}, Bidirectional Gated SSM (BiGS)~\citep{Wang:22}, the H3 model~\citep{Fu:23}, as well as Mamba~\citep{Gu:24arxiv}.

\paragraph{Recurrent Neural Networks.} Recurrent Neural Networks (RNNs) have been proposed as alternatives to Transformer architectures and attention mechanisms, primarily because their computation scales linearly with respect to context length. Deep Linear Recurrent Units (LRUs) within RNNs have demonstrated strong capabilities in language modeling~\citep{Orvieto:23, De:24arxiv}, along with Hierarchically Gated Linear RNNs (HGRN)~\citep{Qin:23} and its successor HGRN2~\citep{Qin:24arxiv}. Among the notable RNN-based models for large-scale language modeling is RWKV~\citep{Peng:23arxivshort,Peng:24arxivshort}, which has shown performance on par with Transformer-based models.

\paragraph{Gating.}
Gating constitutes a fundamental concept originating from LSTM, and it has been independently rediscovered and recontextualized in numerous recent models. Implementations of gating mechanisms are found across various architectures, such as HGRN~\citep{Qin:23}, HGRN2~\citep{Qin:24arxiv}, Gated Linear Attention (GLA)~\citep{Yang:23arxiv}, Gated State Space (GSS) models~\citep{Mehta:22}, Bidirectional Gated SSM (BiGS)~\citep{Wang:22}, Moving Average Equipped Gated Attention (MEGA)~\citep{Ma:22}, RWKV~\citep{Peng:23arxivshort}, and Mamba~\citep{Gu:24arxiv}.

\paragraph{Covariance Update Rule.} 
To improve the memory capacity, we integrate a matrix memory structure into the mLSTM cell, employing a covariance update rule. Variants based on similar update principles include Fast Weight Programmers~\citep{Schmidhuber:92ncfastweights,Schlag:21}, RWKV-5 and RWKV-6~\citep{Peng:24arxivshort}, Retention~\citep{Sun:23arxiv}, Linear Transformer~\citep{Katharopoulos:20}, as well as HGRN2~\citep{Qin:24arxiv}.

\paragraph{Most Related.} From a conceptual standpoint, the models most akin to xLSTM include Retention~\citep{Sun:23arxiv}, RWKV~\citep{Peng:23arxivshort,Peng:24arxivshort}, and HGRN2~\citep{Qin:24arxiv}. These architectures incorporate notions such as matrix memory and/or gating mechanisms. Nonetheless, unlike the newly introduced sLSTM, these methods do not facilitate memory mixing. This capability of memory mixing is crucial for effectively addressing state tracking challenges, which makes LSTMs inherently more expressive than both State Space Models (SSMs) and Transformers~\citep{Merrill:24,Deletang:23}. State tracking plays a vital role in tasks like code execution or maintaining the continuity of entities over extended narratives.

\paragraph{Residually Stacking Architectures.} In line with the architecture of virtually all modern, large-scale deep learning frameworks, xLSTM models are developed by residual stacking of modular components~\citep{Srivastava:15,He:16}.

This approach has been foundational in constructing deep convolutional neural networks~\citep{He:16} and also underpins the design of Transformers~\citep{Vaswani:17}. Transformers serve as the core technological advance driving Large Language Models (LLMs) such as GPT-3~\citep{Brown:20short}, ChatGPT~\citep{Schulman:22short}, GPT-4~\citep{Achiam:23gpt4short}, Megatron-LM~\citep{Shoeybi:19}, Gopher~\citep{Rae:21short}, ERNIE 3.0 Titan~\citep{Wang:21arxivshort}, GLaM~\citep{Du:21glamshort}, Chinese M6~\citep{Lin:21}, multilingual AlexaTM 20B~\citep{Soltan:22}, OPT~\citep{Zhang:22opt}, Chinchilla~\citep{Hoffmann:22short}, BLOOM~\citep{Scao:22arxivshort}, GLM-130B~\citep{Zeng:22arxivshort}, LaMDA~\citep{Thoppilan:22short}, PaLM~\citep{Chowdhery:22short}, Llama~\citep{Touvron:23arxiv}, and Gemini~\citep{GeminiTeam:23arxiv,Reid:24arxiv}.

\section{Experiments}
\label{sec:Exp}

This section presents an empirical analysis of xLSTM, emphasizing its performance in language modeling as compared to established approaches. Section~\ref{sec:ExpSyntethic} explores the unique abilities of xLSTM using controlled synthetic benchmarks. In Section~\ref{sec:ExpComparison}, we measure the validation perplexity across a range of prevailing language modeling methods, all trained on a 15B-token subset of the SlimPajama dataset~\citep{Soboleva:23}. Additionally, we carry out ablation studies for xLSTM on this dataset. The scaling properties of these techniques are further examined, following the methodology of \citet{Kaplan:20} and \citet{Brown:20short}. Section~\ref{sec:ExpLanguage} details a more extensive language modeling evaluation. Here, we train both xLSTM and the leading baselines from Section~\ref{sec:ExpComparison} on a significantly larger set of 300B SlimPajama tokens~\citep{Soboleva:23}. This part comprises four main evaluations: (1) the models’ abilities to generalize to longer input contexts; (2) a comparison based on validation perplexity and downstream task performance following the protocol of~\citep{Sutawika:24}; (3) an assessment over the 571 domains represented in the PALOMA language benchmark~\citep{Magnusson:23arxivshort}; and (4) a renewed scaling study, this time leveraging training runs that are twenty times larger.

\label{sec:xlstm_blocks_notation}
Throughout all experiments, we adopt the notation xLSTM[$a$:$b$], where $a/b$ indicates the proportion of mLSTM-based to sLSTM-based xLSTM blocks. For instance, xLSTM[7:1] specifies a configuration comprising seven mLSTM-based blocks and a single sLSTM-based block, totaling eight blocks. In the case of the commonly used total of 48 blocks, this configuration would include 42 mLSTM-based blocks and 6 sLSTM-based blocks. Additionally, across all experimental setups, we employ up-projection blocks preceding the mLSTM and following the sLSTM blocks, respectively.

\subsection{Synthetic Tasks and Long Range Arena}
\label{sec:ExpSyntethic}
We begin by evaluating xLSTM’s novel exponential gating mechanism, equipped with memory mixing, on tasks involving formal languages~\citep{Deletang:23}. Subsequently, we investigate how the newly introduced matrix memory in xLSTM performs on the Multi-Query Associative Recall benchmark~\citep{Arora:23arxiv}. Lastly, we measure xLSTM’s capacity for handling extended sequences using the Long Range Arena suite~\citep{Tay:21}.

\paragraph{Assessment of xLSTM's Exponential Gating with Memory Mixing.} We evaluate the newly introduced exponential gating combined with memory mixing in xLSTM, designed to tackle state tracking challenges~\citep{Merrill:24, Merrill:23}. To facilitate testing for length extrapolation, we adapt and build upon the formal language tasks originally presented by \citet{Deletang:23}, incorporating support for multi-length training. Appendix~\ref{sec:appExpSynthetic-formalLang} offers an in-depth explanation of each task as well as comprehensive results. 
Our evaluation contrasts xLSTM with alternative architectures, including Transformers, State Space Models, and standard Recurrent Neural Networks. For all methods, we compute accuracy only on tokens central to the defined task, rescaling scores from 0 (indicating random guessing) to 1 (reflecting flawless performance). Specifically, we focus on two-block configurations across various models: xLSTM[0:1] (representing pure sLSTM), xLSTM[1:0] (representing pure mLSTM), xLSTM[1:1], Llama, Mamba, RWKV, Retention, Hyena, LSTM, and LSTM used within Transformer block structures (denoted LSTM (Block)). These comparative outcomes appear in Figure~\ref{fig:formal-main}.
Architectures such as Transformers or State Space Models, which lack memory mixing and therefore cannot perform state tracking, fail to handle certain regular languages—for instance, the parity task. This observation is consistent with previous work that demonstrates a fundamental expressive shortcoming of Transformers and State Space Models relative to RNNs~\citep{Merrill:24, Merrill:23, Deletang:23}.

\paragraph{Test of xLSTM's Memory Capacities on Associative Recall Tasks.} This study evaluates the matrix-based memory module introduced in xLSTM by measuring its performance on the Multi-Query Associative Recall task~\citep{Arora:23arxiv}. In this setup, each input sequence includes a set of randomly sampled key-value pairs from a comprehensive vocabulary, which the model must remember for subsequent querying. To further challenge the models, we expand upon the baseline task by raising the number of key-value pairs to as many as 256, while lengthening the input sequence to 2048 tokens. This extension enables a more thorough investigation of the memory limits across a variety of models. Specifically, we benchmark the performance of two-block configurations for Llama, Mamba, RWKV-5, RWKV-6, as well as xLSTM[1:1] and xLSTM[1:0]. Evaluation is performed based on the correctness of retrieving the target pairs. Since the memory capacity of Transformers (such as Llama) scales exponentially with the dimension of encoding~\citep{Ramsauer:21}, they serve as the reference point for performance in this benchmark. Figure~\ref{fig:mqar-main} summarizes the findings. 
xLSTM[1:1] emerges as the top-performing non-Transformer model, including in smaller-scale settings. Notably, the inclusion of the sLSTM block does not impede recall ability; instead, it enhances it—an effect that becomes most pronounced on tasks involving 256 key-value pairs. Appendix~\ref{sec:appExpSynthetic-mqar} contains further analyses, where extrapolation experiments suggest that xLSTM’s improved memory design allows the model to generalize effectively to sequence lengths that extend beyond those present in its training data.

\paragraph{Test of xLSTM's Long Context Capabilities on Long Range Arena.} To evaluate how well xLSTM processes extended sequences and broad contextual information, we benchmark various approaches using the Long Range Arena~\citep{Tay:21}. Across all evaluated tasks, xLSTM consistently yields high performance, indicating that its architecture is highly effective at managing the challenges inherent in long-context scenarios.

For more details, see Appendix~\ref{sec:appExpSynthetic-lra}.

\subsection{Method Comparison and Ablation Study}
\label{sec:ExpComparison}

In order to investigate the central research question—namely, the capabilities of our newly introduced LSTM variants in large-scale language modeling—we conduct experiments where xLSTMs, Transformers, State Space Models, and additional approaches are trained on 15B tokens sourced from SlimPajama, employing a consistent auto-regressive framework. We then evaluate these models on a shared validation set and carry out ablation analyses specifically for the xLSTMs.

\paragraph{Comparative Analysis of xLSTM and Baseline Techniques.} For benchmarking purposes, we trained all models on a corpus of 15B tokens sourced from SlimPajama~\citep{Soboleva:23}. The main metric for evaluation is validation set perplexity. We consider the following competitors: xLSTM (the proposed approach), GPT-3 (Transformer-based)~\citep{Brown:20short}, Llama (Transformer-based)~\citep{Touvron:23arxiv}, H3 (based on State Space Model, SSM)~\citep{Fu:23}, Mamba (SSM)~\citep{Gu:24arxiv}, three versions of RWKV (RWKV-4, RWKV-5, RWKV-6; RNN-based)~\citep{Peng:23arxivshort,Peng:24arxivshort}, GLA (a linear Transformer variant)~\citep{Yang:23arxiv}, HGRN and HGRN2 (RNN-based)~\citep{Qin:23,Qin:24arxiv}, as well as RetNet~\citep{Sun:23arxiv} and Hyena~\citep{Poli:23} (linear Transformer models). We also include specific xLSTM configurations: xLSTM[1:0] and xLSTM[7:1], as described in Section~\ref{sec:xlstm_blocks_notation}. 
All training employed mixed precision. However, for RWKV-5, RWKV-6, GLA, and HGRN2, mixed-precision training did not leverage PyTorch’s automated mixed precision functionality (refer to Appendix Section~\ref{sec:appGenTrainProc} for additional details).
Methods under study are classified into three categories: (a) Transformers, (b) State Space Models (SSMs), and (c) Recurrent Neural Networks (RNNs) alongside linear Transformers, which represent linearized variants substituting the attention component of standard Transformers. All models are parameter-matched to GPT-3’s configuration with 350M parameters—specifically, embedding dimension 1024 and 24 residual layers. Notably, among these, GPT-3 uniquely applies weight sharing between token and output embeddings, thus containing a slightly lower parameter count. As evidenced by Table~\ref{tab:spaj15b_model_results}, xLSTM achieves superior validation perplexity across all baselines. The Appendix~\ref{sec:appExpComparison} provides comprehensive comparative results.
Scaling analysis, depicted in Figure~\ref{fig:temp_sclaw15B}, suggests that xLSTM sustains its performance gains for model configurations of larger sizes.

\paragraph{Ablation Studies.}
As presented in Table~\ref{tab:spaj15b_model_results} and Figure~\ref{fig:temp_sclaw15B}, xLSTM demonstrates strong performance in language modeling tasks when trained on a 15B-token subset of SlimPajama. To systematically evaluate the impact of each architectural modification from LSTM to xLSTM, we gradually adapt a standard LSTM architecture toward the xLSTM design. The first modification introduces LSTM layers into pre-LayerNorm residual backbones. Next, we enhance the model by adding a post up-projection block. In the final step, exponential gating and matrix memory mechanisms are incorporated.

The results are shown in Table~\ref{tab:ablstudies} (top). 

Ablation experiments indicate that the marked gains in performance are largely due to the combination of exponential gating and matrix memory. Furthermore, considering the critical role of gating within RNNs and State Space Models, various gating strategies are systematically removed and evaluated. As shown in Table~\ref{tab:ablstudies} (bottom), results suggest that configuring each gate to be trainable and responsive to the input leads to successive improvements. Further detailed analyses concerning the separate backbone modules can be found in Appendix~\ref{sec:appAblDetails}.

\subsection{xLSTM as Large Language Model}
\label{sec:ExpLanguage}

This research concludes with extensive language modeling experiments aimed at evaluating the viability of xLSTM as a large language model (LLM). To this end, we expand the training dataset to 300B tokens sourced from SlimPajama, matching the token count used in models such as Mamba~\citep{Gu:24arxiv} and Griffin~\citep{De:24arxiv}.

xLSTM is benchmarked against RWKV-4, Llama, and Mamba, each representing a different methodological class as outlined in Section~\ref{sec:ExpComparison}. RWKV-4 serves as the RNN exemplar in these comparisons, as robust precision configurations for RWKV-5, RWKV-6, and HGRN2 were not established until after initiation of the 300B-token experiments (see Appendix~\ref{sec:appGenTrainProc}).

We implement multiple model sizes (125M, 350M, 760M, 1.3B), evaluating each for their ability to extrapolate across varying sequence lengths and measuring their performance on a designated validation set. Additionally, we examine outcomes on downstream tasks, analyze results in language modeling across 571 distinct domains using the PALOMA benchmark, and further explore each model’s scaling law dynamics.

\paragraph{Sequence Length Extrapolation.}
\label{sec:spaj300B_ctx_extrapolation}
To begin, we examine the capacity for sequence length extrapolation in 1.3B parameter models, specifically comparing xLSTM, RWKV-4, Llama, and Mamba. Each model is trained with a fixed context length of 2048 tokens and subsequently assessed on extended context lengths up to 16384 tokens. The corresponding outcomes are depicted in Figure~\ref{fig:spaj300B_seq_extrapolation}. Notably, xLSTM exhibits sustained low perplexity across these extended contexts, setting it apart from the other architectures evaluated.

\paragraph{Validation Perplexity and Downstream Tasks.}
\label{sec:spaj_downstream_eval}
Next, we assess models of various sizes—including xLSTM, RWKV-4, Llama, and Mamba—on the SlimPajama validation set by measuring next token prediction accuracy. We further test these models on a range of downstream benchmarks designed to evaluate performance in common sense reasoning tasks.

The third column in Table~\ref{tab:lmeval} presents the validation set perplexities for the various methods compared. Notably, xLSTM[1:0] and xLSTM[7:1] achieve the lowest perplexity scores across all model sizes on the validation set. Performance outcomes for downstream tasks are shown in the remaining columns of Table~\ref{tab:lmeval}. In nearly all tasks and model configurations, xLSTM consistently outperforms other approaches, with the exception of certain instances on the ARC benchmark where Mamba attains the top results. For additional information, refer to Appendix~\ref{sec:appExpLanguage}.

\paragraph{Performance on PALOMA Language Tasks.} Third, we evaluate next token prediction across all model sizes for xLSTM, RWKV-4, Llama, and Mamba, using the PALOMA language tasks~\citep{Magnusson:23arxivshort}. Perplexity is reported as a measure of next token prediction capability across 571 textual domains, with coverage spanning from sources such as nytimes.com to the Reddit community r/depression.

Results are summarized in Table~\ref{tab:ppleval}, which separates perplexity scores for language modeling benchmarks (the first seven columns) from those for specific domain tasks (the final five columns). On these benchmarks, xLSTM[1:0] achieves lower perplexity values compared to xLSTM[7:1].

More specifically, xLSTM[1:0] outperforms Mamba in 568 out of 571 (99.5\%) domains, Llama in 486 out of 571 (85.1\%) domains, and RWKV-4 in 570 out of 571 (99.8\%) domains. Additional results can be found in Appendix~\ref{sec:appExpLanguage}.

\paragraph{Scaling Laws.} Fourth, we investigate power-law scaling properties, which enable predictions of model performance as size increases~\citep{Kaplan:20,Brown:20short}. Figure~\ref{fig:temp_sclaw300B} illustrates these scaling trends. While all models demonstrate comparable scaling patterns, they differ in baseline performance levels. Among them, RWKV-4 exhibits the least favorable results, followed in succession by Llama and then Mamba. xLSTM outperforms Mamba, with the degree of improvement over Mamba being comparable to Mamba's margin over Llama. The observed scaling relationships suggest that as models become larger, xLSTM is likely to maintain or extend its advantage relative to both Transformer architectures and State-Space models.

\textbf{Generation Times and Maximal Throughput.} In Figure~\ref{fig:inference_time_generation}, we present measurements of text generation latency for the various xLSTM configurations at the 1.3B parameter level. Maximal throughput results are displayed in Figure~\ref{fig:inference_time_decoding_throughput} (left). These xLSTM results are evaluated alongside equivalently-sized implementations of Mamba, Llama, and RWKV available from HuggingFace, with the Llama model utilizing a static key-value cache. At the time these experiments were conducted, neither compilation with a complete cache for the Transformer Model nor the use of \texttt{torch.compile} with the Mamba model were operational. All text generation measurements employ a batch size of 1 and a pre-fill of 16 tokens across the models, which is anticipated to be the most advantageous setup for the Transformer architecture.

The results depicted in Figure~\ref{fig:inference_time_generation} illustrate that xLSTM, along with the recurrent architectures Mamba and RWKV-4, demonstrates linear scaling in generation time, as opposed to the quadratic complexity exhibited by Llama. In the decoding throughput evaluations, we report the effects of varying both the batch size and the pre-fill parameter specifically for the Llama model. The data in Figure~\ref{fig:inference_time_decoding_throughput} (right) reveal that xLSTM’s memory requirements do not grow with batch size, enabling substantially larger batches and resulting in superior maximum throughput relative to Llama.

\section{Limitations}

(i) Unlike mLSTM, the memory mixing approach of sLSTM inherently prevents parallelizable operations, precluding a highly parallel implementation. Despite this, we have implemented a custom CUDA kernel for sLSTM that achieves execution speed less than twice that of our parallelized mLSTM version. (ii) The current mLSTM CUDA kernels are not specifically tuned for performance and thus the existing version operates at roughly a quarter of the speed of FlashAttention or the scan approach used in Mamba. Enhanced CUDA kernel efficiency, similar to FlashAttention, could likely be achieved with targeted optimization. (iii) Because mLSTM relies on matrix memory, the computational demand is elevated, as processing involves $d \times d$ matrices. Nevertheless, both memory update and read operations remain parameter-free and fully parallelizable using standard matrix operations, so the increased memory complexity introduces only minimal wall time overhead. (iv) Careful selection of initial values for the forget gates is essential. (v) Given that the matrix memory does not depend on sequence length, expanding the sequence can potentially exceed the memory’s capacity for long contexts. However, empirical evidence suggests this is not a significant issue for sequences up to 16k, as discussed in Section~\ref{sec:spaj300B_ctx_extrapolation}. (vi) Owing to the heavy computational resources required for large-scale language model experiments, neither architectural nor hyperparameter optimization was performed exhaustively for xLSTM, particularly at larger scales. We expect that comprehensive optimization will be critical for xLSTM to achieve optimal performance.

\section{Conclusion}
This work addresses the core inquiry: How effectively does language modeling perform when LSTM architectures are scaled to the order of billions of parameters? Our findings indicate that the capabilities extend “at least as far as” those offered by prevailing models such as Transformers and State Space Models. We present an enhanced variant, xLSTM, featuring exponential gating combined with memory mixing and an updated memory structure. In empirical evaluations, xLSTM shows strong performance in language modeling, outperforming or matching leading models including Transformers and State Space Models. Moreover, scaling behaviors observed suggest that expanding the parameter count in xLSTM will make it a formidable rival to the Large Language Models constructed via Transformer architectures. Beyond language modeling, the innovations in xLSTM hold significant promise for advancing research in areas such as Reinforcement Learning, Time Series Prediction, and the modeling of physical processes.

\end{document}